%==============================================================================
% MATH 5080
% Class:        10
% Date:         Sept 29th
% cross checked against lecture notes
%==============================================================================

\subsection{Properties of continuous functions: compactness}

\begin{remark}
    Continuity does not preserve closedness.
\end{remark}
\begin{example}
    Consider $D = [1, \infty), f(x) = \frac{1}{x}$. Then the domain is closed,
    but the image $f(D) = (0, 1]$ and is not closed.
\end{example}
\begin{remark}
    Continuity does not preserve boundedness.
\end{remark}
\begin{example}
    Consider $D = [0, 1], f(x) = \frac{1}{x}$. 
    Then the domain is bounded,
    but the image $f(D) = [1, \infty)$ and is not bounded.
\end{example}

\begin{theorem}[Continuity preserves compactness]
    Let $D \subset \mathbb{R}$ be compact, and $f: D \to \mathbb{R}$ be continuous.
    Then, $f(D)$ is compact.
\end{theorem}

\begin{proof}
    Let $\mathcal{F} = \{ O_{\alpha} \}_{\alpha \in \Delta} $
    be an open cover for $f(D)$. WTS there exists
    a finite subcover for $f(D)$. 

    \medskip

    By theorem, each $\forall O_{\alpha}$, $\exists H_{\alpha}$
    open s.t. $f^{-1} (O_{\alpha}) = H_{\alpha} \cap D$.
    \medskip

    Fix $x \in D$, $f(x) \in f(D) \implies$
    $\exists \alpha \in \Delta$ s.t. 
    $f(x) \in O_{\alpha}$ 
    $\implies$
    $x \in f^{-1}(O_{\alpha}) = H_{\alpha} \cap D \subset H_{\alpha}$.
    Hence $ \{ H_{\alpha} \}_{\alpha \in \Delta} $ is an open cover
    for $D$.
    \medskip

    Since $D$ compact, $\exists \alpha_1, \alpha_2, \hdots \alpha_k$
    s.t. 
    $D \subset H_{\alpha_1} \cup \hdots \cup H_{\alpha_k}$.
    \medskip

    Fix $y \in f(D), \exists x \in D$ s.t. 
    $y = f(x)$. Since $x \in H_{\alpha_i}$
    for some $i = 1, \hdots k$, we have 
    $x \in H_{\alpha_i} \cap D = f^{-1}(O_{\alpha_i})$
    $\implies$
    $y = f(x) \in O_{\alpha_i}$.
\end{proof}


\begin{corollary}[Extreme Value Theorem]
    Let $D$ be compact, and $f: D \to \mathbb{R}$ be continuous. 
    Then, $f$ attains both its maximum and minimnum on $D$, i.e.
    $\exists x_1, x_2 \in D$ s.t. $f(x_1) \leq f(x) \leq f(x_2)$ 
    $, \forall x \in D$.
\end{corollary}

\begin{proof}
    By theorem, $f(D)$ is compact.
    \medskip

    \textbf{Claim}: Suppose $S \subset \mathbb{R}, S \neq \emptyset$
    and $S$ compact, then $S$ contains its max and min. 

    \textbf{Proof}: since $S$ bounded and nonempty $\sup S, \inf S$ exists.
    Since $S$ closed, $\inf S, \sup S \in S$. 
\end{proof}

\begin{definition}[Uniform continuity]
    Let $f: D \to \mathbb{R}$, we say $f$ is uniformly continuous on $D$ if 
    $\forall \ep > 0, \exists \delta > 0$ s.t. $\forall x, y \in D$, 
    $|x - y| < \delta \implies |f(x)- f(y) | < \ep$.
\end{definition}

\begin{remark}
    $f$ can be continuous at any $c \in D$ but not uniformly continuous.
\end{remark}

\begin{example}
    $f(x) = \frac{1}{x}, D = \left( 0, 1 \right) $ is continuous but 
    not uniformly continuous.
\end{example}

\begin{example}
    $f(x) = \frac{1}{x}, D = (1, 2]$ is uniformly contiuous. 
    Fix $\ep > 0$, we seek $\delta > 0$ such that $|x-y| < \delta \implies |f(x) - f(y) | < \epsilon$. 
    \[
        |f(x) - f(y)|= \left| \frac{1}{x} - \frac{1}{y} \right|  = 
        \frac{ \left| x- y \right| }{x-y} \leq |x-y|
    \]
    Hence we can pick $\delta = \ep$.
\end{example}

\begin{example}
    $f(x) = \sin x$ is uniform continuous on $\mathbb{R}$.
\end{example}

\begin{example}
    $f(x) = x$ is uniform continuous on $\mathbb{R}$,
    $g(x) = x^{2}$ is not.
\end{example}


\begin{remark}
    Boundedness is not required for uniform continuity. Consider $f(x) = x$.
\end{remark}

\begin{theorem}
    Let $D$ be compact and $f: D \to \mathbb{R}$ be continuous on $D$. Then $f$ is uniformly
    continuous on $D$.
\end{theorem}

\begin{remark}
    On compact sets, continuity is necessary and sufficient for uniform continuity.
\end{remark}

\begin{proof}
    Fix $ \ep > 0$. $\forall x \in D$, since $f$ continuous at $x$, $\exists \delta_x$ s.t. 
    \[
        |x-y| < \delta_x \implies \left| f(x) - f(y) \right|  < \frac{\ep}{2}
    \]
    Fix $y_1, y_2 \in N(x, \delta_x)$. We want to say that $|f(y_1) - f(y_2) | $ is small enough.
    \medskip

    Define $O_{x} = N \left( x, \frac{\delta_x}{2} \right) $. Then $ \{ O_x \}_{x \in D} $ is an open cover for $D$. 
    By compactness, $\exists x_1, x_2, \hdots x_n \in D$ \text{ s.t. } $D \subset O_{x_1} \cup \hdots O_{x_k}$.
    Define 
    \[
        \delta \coloneq \min \{ 
            \frac{\delta_{x_1}}{2},
            \frac{\delta_{x_2}}{2},
            \hdots,
            \frac{\delta_{x_k}}{2},
        \} 
    \]
    Claim: $\forall x, y \in D, |x-y| < \delta \implies |f(x) - f(y)| < \ep$. 
    \medskip

    Fix $x, y$ with $|x - y| < \delta$.
    \begin{align*}
        & |x-y| < \delta, \exists O_{x_i} \text{ s.t. } x \in O_{x_i} \\
        \implies & \left| f(x) - f(x_i) \right|  < \frac{\ep}{2} \\
    \end{align*}
    Also
    \begin{align*}
        & |x-y| < \delta \leq \frac{\delta_{x_i}}{2} \\
        \implies & |y - x_i| \leq |y - x| + |x - x_i| < 
        \frac{\delta_{x_i}}{2} + 
        \frac{\delta_{x_i}}{2} 
        = \delta_{x_i} \\
        \implies & y \in N(x_i, \delta_{x_i}) \\
        \implies & \left| f(y) - f(x_i) \right| < \frac{\ep}{2}
    \end{align*}
    Hence
    \begin{align*}
        |f(x) - f(y)| \leq
        \left| f(x) - f(x_i) \right|  + 
        \left| f(x_i) - f(y) \right| 
        <
        \frac{\ep}{2} + \frac{\ep}{2} = \ep
    \end{align*}
\end{proof}

\begin{definition}
    We say $\tilde{f}: E \to \mathbb{R}$ is an extension 
    of a function $f: D \to \mathbb{R}$ if $D \subset E$
    and $\tilde{f}(x) = f(x), \forall x \in D$.
\end{definition}


\begin{theorem}
    \label{thm:equivalent-characterization-of-uniform-continuity}
    Let $f(a, b) \to \mathbb{R}$. TFAE
    \begin{enumerate}
        \item $f$ is uniformly continuous on $(a, b)$
        \item $f$ is continuous on $(a,b)$ and both $\lim_{x \to a^+}f(x)$ and $\lim_{x \to b^{-1}}f(x)$
        exist in $\mathbb{R}$
        \item $f$ can be extended to a continuous function $\tilde{f}$ on $[a,b]$
    \end{enumerate} 
\end{theorem}

\begin{lemma}[Uniform continuity preserves Cauchyness]
    \label{lemma:uniform-cty-preserves-cauchyness}
    Let $f: D \to \mathbb{R}$ be a uniformly continuous function on $D$, 
    
    then for any Cauchy sequence
    $\{ a_n \} \subset D$, 
    $ \{ f(a_n) \} $ is also Cauchy. 
\end{lemma}

\begin{proof}[Proof of \ref{lemma:uniform-cty-preserves-cauchyness}]
    Let $ \{ a_n \} $ be Cauchy. WTS that $ \{ f(a_n) \} $ is Cauchy. 
    $\forall \ep > 0$, we seek $N > 0 \text{ s.t. } |f(a_n) - f(a_m)| < \ep, \forall n, m > N$. 
    \medskip

    Since $f$ is uniform contiuous on $D$, $\exists \delta \text{ s.t. } |x-y| < \delta \implies |f(x) - f(y) | < \ep$. 
    \medskip

    Since $ \{ a_n \} $ is Cauchy, $\exists N > 0$.s.t 
    $\forall n, m > N$,
    \[
        |a_n - a_m| < \delta \implies |f(a_n) - f(a_m) | < \ep
    \]
\end{proof}

\begin{proof}[Proof of
    \ref{thm:equivalent-characterization-of-uniform-continuity}
]

    \begin{enumerate}
        \item (2) $\implies$ (3) \remind{On homework, revisit}
        \item (3) $\implies$ (2) 
        \begin{itemize}
            \item Fix any $c \in (a, b)$. WTS $f$ continuous at $c$.
            \medskip
            Since $c$ interior, $\exists \delta$ s.t. $N(c, \delta) \subset (a,b)$.
            Hence $f(x) = \tilde{f}(x)$ $\forall x \in N(c, \delta)$.
            $\tilde{f}$ continuous on $[a,b] \implies \tilde{f}$ continuous on $N(c, \delta)$
            $\implies f$ continuous on $N(c, \delta)$. Hence $f$ continuous at $c$.
            \item WTS $\lim_{x \to a^+}f(x)$ exists. 
            Since $\tilde{f}$ continuous at 
            $a$, $\lim_{x \to a^+}f(x) = \lim_{x \to a^+}\tilde{f}(x) = f(a)$.
            \item WTS $\lim_{x \to b^-}f(x)$ exists.
            Since $\tilde{f}$ continuous at 
            $b$, $\lim_{x \to b^-}f(x) = \lim_{x \to b^-}\tilde{f}(x) = f(b)$.
        \end{itemize} 
        \item (3) $\implies$ (1) \remind{On homework, revisit}
        \item (1) $\implies$ (2): it suffices to show that 
        $f$ uniformly continuous on $(a,b)$ implies 
        $\lim_{x \to a^+}f(x)$,
        $\lim_{x \to b^-}f(x)$ exist in $\mathbb{R}$.
        By symmetry, it suffices to show that 
        $\lim_{x \to a^+}f(x)$ exists.
        \medskip 

        ATC otherwise, i.e. $\lim_{x \to a^+}f(x)$ does not exist,
        then by theorem, $\exists \{ a_n \} \subset (a,b)$
        s.t. $a_n \to a^+$
        and $\{ f(a_n) \} $ diverges. 
        However $\{ a_n \} $ is Cauchy which means 
        $\{ f(a_n) \} $ must be Cauhcy. Hence a contradiction.
    \end{enumerate}
\end{proof}

% =================================================================================

\subsection{
    Properties of continuous functions: connectedness and monotonicity
}

\begin{theorem}[Intermediate Value Theorem]
    \label{thm:IVT}
    Let $f$ continuous, $f: [a,b] \to \mathbb{R}$, and $y$ is a value 
    in between $f(a), f(b)$.
    Then, $\exists c \in (a, b) \text{ s.t. } f(c) = y$.
\end{theorem}

% \begin{remark}
%     \textit{If $D$ is connected, then $f(D)$ is connected.}
% \end{remark}

\begin{theorem}[Continuity preserves connectedness]
    \label{thm:cty-preserves-connectedness}
    Let $f: D \to \mathbb{R}$ continuous and $D$ is connected, 
    then $f(D)$ is connected.
\end{theorem}

\begin{corollary}
    Let $I$ be a compact interval (i.e. closed and bounded $[a,b]$
    or a singleton
    ). If $f$ is continuous on $I$, then $f(I)$ is 
    a compact interval.
\end{corollary}


\begin{proof}[Proof of \ref{thm:cty-preserves-connectedness} $\implies$ IVT]
    Follows immediately by definition.
    Since $[a,b]$ is connected, 
    $f([a,b])$ is connected.
    WLOG say $f(a) < f(b)$.
    For any $f(a) < y < f(b)$, $y \in f([a,b])$
    $\implies \exists c \in [a, b]$ s.t. 
    $f(c) = y$.
\end{proof}

\begin{proof}[Proof of \ref{thm:IVT}]
    ATC $f(D)$ disconnected, 
    $\exists A, B$ open s.t. 
    $f(D) \subset A \cup B, A \cap f(D) \neq \emptyset$, 
    $B \cap f(D) \neq \emptyset, A \cap B = \emptyset$.

    Since $f $ continuous, and $A$ open, $\exists$
    open $U$ s.t. $f^{-1}(A) = U \cap D$. 
    Similarly, $\exists$ open $V$ s.t. $f^{-1}(B) = V \cap D$. 
    We then claim 
    \begin{enumerate}
        \item $D \subset U \cup V$
        \item $U \cap V \cap D = \emptyset$
        \item $U \cap D \neq \emptyset, V \cap D \neq \emptyset$
    \end{enumerate} 
    Hence $D$ is disconnected.
\end{proof}
\begin{remark}
    Note that (2) above does not necessarily
    imply $U \cap V = \emptyset$. 
    \remind{
        Refer to notes on connectedness which shows that (2)
        is sufficient for disconnectedness.
    }
\end{remark}





























